\documentclass{article}
\usepackage{amsmath}
\usepackage{cancel}
\usepackage{graphicx} % for creating figures

\title{Calculus of Variations}

\begin{document}
\maketitle
So far we have seen the utility of d'Alembert's principle and the Euler-Lagrange
equations derived from it. We would now like to present a new formulation of
dynamics with a new principle that will allow us to re-derive the Euler-Lagrange
equations while providing an even more powerful framework for describing
physical laws.

The mathematics behind this formalism is called the \textit{Calculus of
  Variations} and is concerned with finding the extreme values of
functions-of-functions known as \textit{functionals}. For example, imagine we
have some function $y$ of the variable $x$. Then, our experience with
Lagrangians suggests we should care about objects of the form
\begin{equation}
  \phi\Big( y(x),\; y'(x),\; x\Big)
\end{equation}
that is, for a given value of $x$, $\phi$ maps the function $y$ and its
derivative and maps them to some number. Integrals tend to do this operation of
sending functions to numbers nicely. Therefore, let's consider a scenario such
as
\begin{equation}
  \int\limits_{x_1}^{x_2}\phi\Big( y(x),\; y'(x),\; x\Big)dx
\end{equation}
which looks more familiar if we think of the problem of finding the curve which
minimizes the distance between two points.  We know distance is calculated by 
\begin{equation}
  ds = \Big(dx^2+dy^2 \Big)^{1/2} = \left(1+\left(\frac{dy}{dx}\right)^2\right)^{1/2}dx
\end{equation}
so that the arclength between two points on a curve given by $y(x)$ is
\begin{equation}
  \int\limits_{x_1}^{x_2}\left( 1+y'^2 \right)^{1/2}dx
\end{equation}
To solve the problem, we consider an infinitesimal displacement of the path
given by $y(x)$
\begin{equation}
  Y(x) = y(x)+\varepsilon\eta(x)
\end{equation}
where the two paths share the same endpoints and $\varepsilon$ is a small amount.  

\begin{center}
  \centering
  \includegraphics{./figures/calculus-of-variations-1.svg}
\end{center}

Certainly, in order for the curves to share the same endpoints, we must have
\begin{equation}
  \eta(x_1) = \eta(x_2) = 0 
\end{equation}

If we call our integral $S$ (for sum) then we now have turned our
\textit{functional} into \textit{a function of $\varepsilon$}.
\begin{equation}
  S(\varepsilon) = \int\limits_{x_1}^{x_2}\phi\left(y(x)+\varepsilon\eta(x),\;
    y'(x)+\varepsilon\eta'(x),\; x  \right)dx
\end{equation}

Since we are working with infinitesimal quantities, let's Taylor expand $\phi$
around the original curve (i.e. $\varepsilon = 0 $)
\begin{equation}
  \phi(Y, Y', x) \approx 
\end{equation}


\end{document}






























